% Standalone problem document, generated from the adjacent README.md.
% Compile with XeLaTeX. Mathematical content is maintained in Markdown.
\documentclass[11pt,a4paper]{article}
\usepackage[fontsize=10.5pt]{scrextend}
\usepackage[margin=22mm,headheight=15pt,headsep=9mm,footskip=11mm]{geometry}
\usepackage{fontspec}
\setmainfont{texgyrepagella}[Extension=.otf,UprightFont=*-regular,BoldFont=*-bold,ItalicFont=*-italic,BoldItalicFont=*-bolditalic]
\setsansfont{texgyreheros}[Extension=.otf,UprightFont=*-regular,BoldFont=*-bold,ItalicFont=*-italic,BoldItalicFont=*-bolditalic]
\setmonofont{texgyrecursor}[Extension=.otf,UprightFont=*-regular,BoldFont=*-bold,ItalicFont=*-italic,BoldItalicFont=*-bolditalic,Scale=MatchLowercase]
\usepackage{amsmath,amssymb}
\usepackage{unicode-math}
\setmathfont{texgyrepagella-math.otf}
\usepackage{xcolor,microtype,enumitem,needspace,fancyhdr}
\usepackage{bookmark}
\definecolor{ink}{HTML}{17344B}
\definecolor{accent}{HTML}{197D80}
\definecolor{muted}{HTML}{596773}
\hypersetup{colorlinks=true,linkcolor=accent,urlcolor=accent,pdftitle={MI-15 - A
sum-of-squares representation for the Toeplitz commutator
form},pdfauthor={Open Problems in Numerical Linear Algebra}}
\urlstyle{same}
\setlength{\parindent}{0pt}
\setlength{\parskip}{5pt plus 1pt minus 1pt}
\linespread{1.04}
\setlength{\emergencystretch}{3em}
\widowpenalty=10000
\clubpenalty=10000
\displaywidowpenalty=10000
\allowdisplaybreaks[1]
\setlist{nosep,leftmargin=1.5em}
\providecommand{\tightlist}{\setlength{\itemsep}{0pt}\setlength{\parskip}{0pt}}
\providecommand{\pandocbounded}[1]{#1}
\pagestyle{fancy}
\fancyhf{}
\fancyhead[L]{\sffamily\footnotesize\color{muted}OPEN PROBLEMS IN NUMERICAL LINEAR ALGEBRA}
\fancyhead[R]{\sffamily\bfseries\footnotesize\color{accent}MI-15}
\fancyfoot[L]{\parbox[b]{0.9\textwidth}{\sffamily\fontsize{8}{10}\selectfont\color{muted}Editorial ratings; a bounded search cannot certify that a problem remains unsolved.\\Literature check: 2026-09-12}}
\fancyfoot[R]{\sffamily\footnotesize\color{muted}\thepage}
\renewcommand{\headrulewidth}{0.3pt}
\renewcommand{\headrule}{\hbox to\headwidth{\color{accent}\leaders\hrule height \headrulewidth\hfill}}
\makeatletter
\renewcommand{\section}{\Needspace{5\baselineskip}\@startsection{section}{1}{0pt}{12pt plus 2pt minus 2pt}{4pt}{\sffamily\bfseries\large\color{ink}}}
\renewcommand{\subsection}{\Needspace{4\baselineskip}\@startsection{subsection}{2}{0pt}{9pt plus 1pt minus 1pt}{3pt}{\sffamily\bfseries\normalsize\color{ink}}}
\makeatother
\setcounter{secnumdepth}{0}
\begin{document}
{\sffamily\small\color{muted}Matrix inequalities and norms\par}
\vspace{3pt}
{\raggedright\hyphenpenalty=10000\exhyphenpenalty=10000\sffamily\bfseries\fontsize{21}{24}\selectfont\color{ink}A
sum-of-squares representation for the Toeplitz commutator form\par}
\vspace{7pt}
{\sffamily\small\textbf{Difficulty:} challenging\quad\textbf{Importance:} interesting
to specialist\par}
{\sffamily\small\bfseries\color{accent}Status: Partially resolved\par}
\vspace{4pt}
{\color{accent}\hrule height 0.8pt}
\vspace{6pt}
\textbf{Provenance:} explicit conjecture

\textbf{Rating rationale:} All-order algebraic certification needs more
than known nonnegativity; its direct application is a specific Toeplitz
commutator form.

\subsection{Finite-order SOS certificates --- 12 September
2026}\label{finite-order-sos-certificates-12-september-2026}

\textbf{Author:} Sidney Holden, Center for Computational Biology,
Flatiron Institute, Simons Foundation.
\href{https://github.com/ajt60gaibb/OpenProblemsInNLA/blob/main/references/holden-matrix-2026-09-12/README.md}{Submission
and verified affiliation}.

Sections 1--5 of the
\href{https://github.com/ajt60gaibb/OpenProblemsInNLA/blob/main/references/holden-matrix-2026-09-12/MI-15/proof.md}{proof}
establish sum-of-squares representations for orders \(n=8,9,10,11,12\),
with at most \(2(n-1)^2\) homogeneous quadratic squares. The five
rational Gram certificates passed exact polynomial-identity and integer
positive-definiteness checks. The assertion for every order remains
open.

A separate
\href{https://github.com/ajt60gaibb/OpenProblemsInNLA/blob/main/references/holden-matrix-2026-09-12/verification/MI-15-review.md}{independent
Codex AI-agent audit} passed this limited scope. This is informal
automated review; no complete resolution, historical novelty, external
human peer review or formal verification is claimed. No Lean
verification was performed.

\subsection{Problem statement}\label{problem-statement}

For an integer \(n\ge2\), let
\(x_{-(n-1)},\ldots,x_{n-1},y_{-(n-1)},\ldots,y_{n-1}\) be independent
real variables. Form real Toeplitz matrices \(X=(x_{i-j})_{i,j=1}^n\)
and \(Y=(y_{i-j})_{i,j=1}^n\), and the homogeneous quartic polynomial

\[F_n(x,y)=2\|X\|_F^2\|Y\|_F^2
-2\bigl(\mathop{\mathrm{tr}}\nolimits(X^TY)\bigr)^2-\|XY-YX\|_F^2,\]

where \(\|Z\|_F^2=\sum_{i,j}z_{ij}^2\) for real matrices. Is it true
that for every \(n\ge2\) there are a finite integer \(N_n\ge0\) and
homogeneous quadratic polynomials \(q_{n,1},\ldots,q_{n,N_n}\) with real
coefficients in these \(4n-2\) variables such that

\[F_n(x,y)=\sum_{j=1}^{N_n}q_{n,j}(x,y)^2\]

as a polynomial identity? The polynomials may depend on \(n\). Rational
coefficients are not required.

\subsection{Relevance}\label{relevance}

An explicit sum-of-squares identity would give an algebraic certificate
for a structured matrix inequality and connects Toeplitz calculations to
semidefinite optimization. Nonnegativity alone does not answer this
question.

\subsection{References}\label{references}

\begin{enumerate}
\def\labelenumi{\arabic{enumi}.}
\tightlist
\item
  L. László, \emph{Sum of squares representation for the
  Böttcher--Wenzel biquadratic form}, Acta Universitatis Sapientiae,
  Informatica 4(1) (2012), 17--32: equation (1), §5, and Conjecture 15,
  p.31. \href{https://arxiv.org/pdf/1207.6372}{Primary manuscript}.
\item
  J. Ge, F. Li, Z. Tang, and Y. Zhou, \emph{A survey on the DDVV-type
  inequalities}, Advances in Mathematics (China) 53 (2024), 449--467:
  published Conjecture 4.1, p.461; Conjecture 4.3 in
  \href{https://arxiv.org/html/2402.01085v1}{arXiv:2402.01085v1}.
  \href{https://ccj.pku.edu.cn/Article/DownLoad?id=374327987&type=ArticleFile}{Published
  PDF}.
\end{enumerate}

\subsection{Status check --- 2026-09-10}\label{status-check-2026-09-10}

Both arXiv records remain v1; the published 2024 survey still poses this
SOS assertion. László's particular certificate strategy works through
order seven; its failure in a larger order is not a disproof of all SOS
representations. Searches used \textit{Toeplitz SOS Wenzel},
\textit{Toeplitz Bottcher squares proof}, and
\textit{Toeplitz Conjecture 15 squares}. No general proof or non-SOS
counterexample was located. The already proved Böttcher--Wenzel
nonnegativity statement is not being counted again. The optional
rational-certificate request in the original is not imposed as an
additional conjecture.

\textbf{Audit update (2026-09-10):} Rechecked the survey's Conjecture
4.3 and searched for later Toeplitz SOS results. The general assertion
remains a conjecture; known small-order certificates do not establish
all-order representability. This is a bounded literature check, not a
proof that no solution exists.
\end{document}
